In this section we cover the operations which take collections and return collections.
% \subsection{operation()}
% Consider the expression \inlineocl{src.operation()}, where 
% Formal Definition:
% \begin{equation}\label{def:operation}
%     operation(X,y) \iff \llbracket x_i \neq d \rrbracket
% \end{equation}

% Reformulation using Global Constraints:
% \begin{equation}\label{csp:operation}
%     operation(X,y):
%     \begin{cases}

%     \end{cases}
% \end{equation}
% \newpage
\subsection{select(exp)}
Consider the expression \inlineocl{src.select(exp)}, where \inlineocl{src} is an expression resulting in a collection and \inlineocl{exp} is a predicate: the output of the operation is the sub-collection of elements for which the predicate is true.
This operation can have at most one iterator variable.
% Formal Definition:
% \begin{equation}\label{def:select}
%     select(X,Y,exp) \iff
% \end{equation}

% Reformulation using Global Constraints:
\begin{equation}\label{csp:select}
    select(X,Y,exp):
    \begin{cases}
        y_i = d - (d-x_i)*\llbracket exp(x_i) \rrbracket
    \end{cases}
\end{equation}
For this model we first assume the value is filtered $y_i=d$, then depending on if the expression is true when applied $exp(x_i)$ we add it back the value $x_i$, using $-(d-x_i)$.

% \newpage
\subsection{reject(exp)}
Consider the expression \inlineocl{src.reject(exp)}, where \inlineocl{src} is an expression resulting in a collection of objects or integers and \inlineocl{exp} is a predicate: the output of the operation is the sub-collection of elements for which the predicate is false.
This operation can have at most one iterator variable.
% Formal Definition:
% \begin{equation}\label{def:reject}
%     reject(X,Y,exp) \iff
% \end{equation}

% Reformulation using Global Constraints:
\begin{equation}\label{csp:reject}
    reject(X,Y,exp):
    \begin{cases}
        y_i = d - (d-x_i)*\llbracket \neg exp(x_i) \rrbracket
    \end{cases}
\end{equation}
For this model we first assume the value is filtered $y_i=d$, then depending on if the expression is false when applied $exp(x_i)$ we add it back the value $x_i$, using $-(d-x_i)$.

\newpage
\subsection{any()}
This operation can have at most one iterator variable.
Consider the expression \inlineocl{src.any(exp)}, where \inlineocl{src} is an expression resulting in a collection of objects or integers and \inlineocl{exp} is a predicate: the output of the operation is one of the elements for which the predicate is true.
This operation can have at most one iterator variable.
% Formal Definition:
% \begin{equation}\label{def:any}
%     any(X,Y,exp) \iff
% \end{equation}
% Reformulation using Global Constraints:
\begin{equation}\label{csp:any}
    any(X,y,exp):
    \begin{cases}
        select(X,Y,exp)\\
        asBag(Y,Y')\\
        y = y_1'
    \end{cases}
\end{equation}
For this model, we simply use the selection model to find good values, and use the as bag model to make sure the first value is valid.
Any other way of choosing a valid value could also be chosen.

% \newpage
\subsection{sortedBy(exp)}
This operation can have at most one iterator variable.
Consider the expression \inlineocl{src.sortedBy(exp)}, where \inlineocl{src} is an expression resulting in a collection of objects or integers and \inlineocl{exp} results in an integer: the output of the operation is the sequence of elements with respect to the result of the expression.
% Formal Definition:
% \begin{equation}\label{def:sortedBy}
%     sortedBy(X,Y,exp) \iff
% \end{equation}
Reformulation using Global Constraints:
\begin{equation}\label{csp:sortedBy}
    sortedBy(X,Y,exp):
    \begin{cases}
        stable\_keysort(\langle K,X\rangle,\langle K',Y\rangle,1) \\
        ~~k_i=exp(x_i), \forall i \in [1,z]
        % ~~k'_i=\llbracket y_i=d \rrbracket, \forall i \in [1,z]
    \end{cases}
\end{equation}
For this we simply apply the key sort global constraint, using the result of the expression as key.

% \subsection{collect()}
% Consider the expression \inlineocl{src.collect(exp)}, where exp is an expression.
% Formal Definition:
% \begin{equation}\label{def:collect}
%     collect(X,Y,exp) \iff
% \end{equation}

% Reformulation using Global Constraints:
% \begin{equation}\label{csp:collect}
%     collect(X,Y,exp):
%     \begin{cases}
%         y_i = d - (d-x_i)*\llbracket not exp(x_i) \rrbracket
%     \end{cases}
% \end{equation}